\vspace{-1mm}
\subsection{$E(3)$-Equivariant Neural Networks}
\vspace{-1mm}

$E(3)$-equivariant neural networks (ENNs) are built such that every operation $f$ is equivariant to Euclidean symmetry, satisfying that for $g \in E(3)$, inputs $x$ and scalar weights $w$ and $f:X \times W \rightarrow Y$, $f(D_X(g)x, w) = D_Y(g)f(x, w)$ \cite{tfn,kondor2018clebsch,3dsteerable}.
In addition to the tensor product, equivariant operations also include equivariant activation functions and normalizations as in typical neural networks, and we discuss those used in Equiformer in Sec.~\ref{subsec:equivariant_operations}.


%%%%%%%%%%%%%%%%%%

%Because an ENN is a composition of equivariant operations, it is also equivariant; transformations acting on input spaces propagate to the output spaces, e.g., when we rotate an input atomistic graph, predicted vectors will be rotated.




% \revision{Since each operation is equivariant, the whole networks can propagate transformation acting on input spaces to output spaces and achieve equivariance, e.g., when we rotate an input atomistic graph, predicted vectors will be rotated by design.}
%\revision{There have been several works leveraging equivariant networks~\cite{tfn, 3dsteerable, kondor2018clebsch, l1net, nequip, se3_transformer, segnn, allergo} for different applications such as energy prediction~\cite{tfn, se3_transformer, segnn} and molecular dynamics~\cite{nequip, allergo}.}
%\todo{Typical equivariant operations beyond tensor product (e.g., activation and linear).}
%\todo{How we stack equivariant operations to build equivariant networks.}
%\todo{Some typical examples.}